% =============================================================================
% Module 13: Solutions to Exercises
% =============================================================================
% This file is conditionally included based on \ifsolutions
% =============================================================================

\section{Solutions to Exercises}
\label{sec:micro_solutions}

\noindent
Full worked solutions are also available in the Mathematica script
\solnlink{13_Microlensing/problems\_13.wl}{problems\_13.wl},
which verifies each answer symbolically and numerically.

% ----- Exercise 13.1 -----
\subsection*{Exercise~\ref{ex:micro_magnification}: Total Magnification and Its Limits}

\textbf{(a)} The lens equation $u = y - 1/y$ is $y^2 - u\,y - 1 = 0$, with
roots
\begin{equation*}
    y_\pm = \frac{1}{2}\left(u \pm \sqrt{u^2 + 4}\right) \,.
\end{equation*}
The separation is $y_+ - y_- = \sqrt{u^2+4}$. \checkmark

\textbf{(b)} With $\mu = [1 - y^{-4}]^{-1}$ evaluated at $y_\pm$ one finds
\begin{equation*}
    \mu_\pm = \frac{1}{2}\left[\frac{u^2+2}{u\sqrt{u^2+4}} \pm 1\right] \,,
\end{equation*}
so $\mu_+ > 0$, $\mu_- < 0$, and (adding the absolute values)
\begin{equation*}
    A(u) = |\mu_+| + |\mu_-| = \mu_+ - \mu_-
    = \frac{u^2 + 2}{u\sqrt{u^2 + 4}} \,. \quad \checkmark
\end{equation*}

\textbf{(c)} As $u\to0$, $A(u) \simeq 2/(u\cdot2) = 1/u \to \infty$
(high magnification).  As $u\to\infty$, $A(u) \simeq u^2/(u\cdot u) = 1$
(unlensed). \checkmark

\textbf{(d)} For $u_0 = 1$, $A(1) = 3/\sqrt5 \approx 1.342$, so
\begin{equation*}
    \Delta m = 2.5\log_{10}A(1) \approx 0.32~\text{mag} \,. \quad \checkmark
\end{equation*}


% ----- Exercise 13.2 -----
\subsection*{Exercise~\ref{ex:micro_timescale}: Einstein Radius and Timescale}

\textbf{(a)} With $\Dds = \Ds - \Dd = 4$~kpc, $\Dd = 4$~kpc, $\Ds = 8$~kpc,
$M = 0.3\,\Msun$, eq.~\eqref{eq:micro_thetaE} gives
\begin{equation*}
    \thetaE = \sqrt{\frac{4\Gnewton M}{\speedoflight^2}\,
    \frac{\Dds}{\Dd\,\Ds}} \approx 0.55~\text{mas} \,,
    \qquad
    \RE = \Dd\,\thetaE \approx 2.2~\text{AU} \,.
\end{equation*}

\textbf{(b)} With $v_\perp = 200$~km\,s$^{-1}$,
\begin{equation*}
    t_{\mathrm{E}} = \frac{\RE}{v_\perp}
    = \frac{2.2~\text{AU}}{200~\text{km s}^{-1}}
    \approx 19~\text{days} \,.
\end{equation*}
Galactic events last of order weeks because $\RE\sim$~AU and Galactic
transverse speeds are $\sim\!200$~km\,s$^{-1}$, so $\RE/v_\perp$ is weeks.

\textbf{(c)} $\thetaE \propto \sqrt{M}$ and $\RE = \Dd\thetaE$, so
$t_{\mathrm{E}} = \Dd\thetaE/v_\perp \propto \sqrt{M}\,\Dd^{1/2}\ldots$ ---
more precisely, doubling nothing but the mass scales
$t_{\mathrm{E}} \to \sqrt{M}$: $t_{\mathrm{E}}(4M)/t_{\mathrm{E}}(M) = 2$.
A single measured $t_{\mathrm{E}}$ is one number constraining the three
unknowns $(M, \Dd, v_\perp)$, so it cannot determine any of them alone.  The
degeneracy is broken by finite-source effects, microlensing parallax, or the
astrometric centroid shift (Exercise~\ref{ex:micro_astrometric}). \checkmark


% ----- Exercise 13.3 -----
\subsection*{Exercise~\ref{ex:micro_optical_depth}: Optical Depth and the MACHO Test}

\textbf{(a)} The optical depth is the total Einstein-disk area per unit sky
area, $\tau = n_{\mathrm{ang}}\,\pi\thetaE^2$, where $n_{\mathrm{ang}}$ is the
lens number per steradian.  Each Einstein area scales as
$\pi\thetaE^2 \propto M$, while for fixed mass density the number density
scales as $n \propto 1/M$.  The two factors cancel, leaving $\tau$ dependent
only on the total mass density, not on how it is divided into lenses. \checkmark

\textbf{(b)} For a thin screen at the lens plane, the number of lenses per
steradian is $n_{\mathrm{ang}} = \Sigma\,\Dd^2/M$ and the Einstein disk has
angular area $\pi\thetaE^2$ with $\thetaE^2 = 4\Gnewton M\Dds/
(\speedoflight^2\Dd\Ds)$.  Hence
\begin{equation*}
    \tau = n_{\mathrm{ang}}\,\pi\thetaE^2
    = \frac{\Sigma\,\Dd^2}{M}\cdot
    \frac{4\pi\Gnewton M\,\Dds}{\speedoflight^2\Dd\Ds}
    = \Sigma\cdot\frac{4\pi\Gnewton\,\Dd\,\Dds}{\speedoflight^2\Ds}
    = \frac{\Sigma}{\Sigmacr} \,,
\end{equation*}
using $\Sigmacr = \speedoflight^2\Ds/(4\pi\Gnewton\Dd\Dds)$. \checkmark

\textbf{(c)} With $\tau \sim 5\times10^{-7}$, roughly $1/\tau \sim 2\times10^6$
stars must be monitored to expect one lensed above the $u=1$ threshold at any
instant --- of order a few million. \checkmark

\textbf{(d)} The measured $\tau_{\mathrm{LMC}}\approx1.2\times10^{-7}$ is about
$4$--$5$ times smaller than the value expected if the entire halo were compact
objects.  Combined with the small number of detected events, this limits
MACHOs to $\lesssim25\%$ of the halo mass over the probed mass range, and the
surviving events are largely attributable to ordinary stars in the disk or the
Magellanic Clouds.  Compact baryonic objects are therefore \emph{not} the
dominant dark matter. \checkmark


% ----- Exercise 13.4 -----
\subsection*{Exercise~\ref{ex:micro_peak}: Peak Magnification and Event Width}

\textbf{(a)} Solving $A(u_0) = (u_0^2+2)/(u_0\sqrt{u_0^2+4}) = 5$ numerically
gives $u_0 \approx 0.203$.

\textbf{(b)} For $u_0 = 0.05$, $A(u_0) = 20.02\ldots$ while $1/u_0 = 20$; the
fractional error is $\sim\!0.1\%$, well under $1\%$. \checkmark

\textbf{(c)} The threshold $u = 1$ is reached when
$u_0^2 + (t-t_0)^2/t_{\mathrm{E}}^2 = 1$, i.e.
\begin{equation*}
    |t - t_0| = t_{\mathrm{E}}\sqrt{1 - u_0^2}
    \qquad (u_0 < 1) \,.
\end{equation*}
Because $A(u)$ is monotonic in $u$, $A(u(t)) = A(1)$ if and only if
$u(t) = 1$. \checkmark


% ----- Exercise 13.5 -----
\subsection*{Exercise~\ref{ex:micro_astrometric}: Astrometric Centroid Shift}

\textbf{(a)} Differentiating $\delta(u) = u/(u^2+2)$ (in units of $\thetaE$),
\begin{equation*}
    \frac{d\delta}{du} = \frac{(u^2+2) - u(2u)}{(u^2+2)^2}
    = \frac{2 - u^2}{(u^2+2)^2} = 0
    \quad\Longrightarrow\quad u = \sqrt2 \,.
\end{equation*}

\textbf{(b)} Evaluating at $u = \sqrt2$:
\begin{equation*}
    \delta_{\max} = \frac{\sqrt2}{2 + 2} = \frac{\sqrt2}{4}
    = \frac{1}{2\sqrt2} = 8^{-1/2} \approx 0.354
    \quad(\text{in units of }\thetaE) \,. \quad \checkmark
\end{equation*}

\textbf{(c)} The centroid ellipse has \emph{angular} size $\propto\thetaE
\propto\sqrt{M}$, so measuring it yields $\thetaE$ directly --- unlike the
photometric light curve, which yields only the degenerate combination
$t_{\mathrm{E}} = \Dd\thetaE/v_\perp$.  Combined with the relative parallax
$\pi_{\mathrm{ds}}$, and using $\thetaE^2 = (4\Gnewton/\speedoflight^2)
(M/\mathrm{AU})\,\pi_{\mathrm{ds}}$ expressed in milliarcseconds, the lens mass
is $M = 0.123\,\Msun\,\thetaE^2/\pi_{\mathrm{ds}}$; the coefficient
$1/8.14 \approx 0.123\,\Msun\,\mathrm{mas}^{-2}$ is verified in
\texttt{problems\_13.wl}. \checkmark


% ----- Exercise 13.6 -----
\subsection*{Exercise~\ref{ex:micro_quasar}: Quasar Microlensing Scales}

\textbf{(a)} The Einstein radius projected into the source plane is
$r_{\mathrm{E}} = \sqrt{4\Gnewton M\,\Ds\Dds/(\speedoflight^2\Dd)}$.  For
$M = \Msun$ with representative angular-diameter distances
($\Dd\approx1250$~Mpc at $z_d=0.5$, $\Ds\approx1750$~Mpc at $z_s=2$,
$\Dds\approx1550$~Mpc), $r_{\mathrm{E}} \approx 6\times10^{16}$~cm --- of order
$10^{16}$~cm, consistent with eq.~\eqref{eq:micro_qso_rE}. \checkmark

\textbf{(b)} All mass-independent factors are fixed, so
$r_{\mathrm{E}} \propto \sqrt{M}$; $r_{\mathrm{E}}(4M)/r_{\mathrm{E}}(M) = 2$.
\checkmark

\textbf{(c)} With $R = 10^{15}$~cm and $v_\perp = 600$~km\,s$^{-1}$,
\begin{equation*}
    t_{\mathrm{cross}} = \frac{R}{v_\perp}
    = \frac{10^{15}~\text{cm}}{600~\text{km s}^{-1}}
    \approx 6~\text{months} \,.
\end{equation*}
The full Einstein-crossing time $t_{\mathrm{E}}\sim15$~yr is impractically
long, but each \emph{caustic} crossing produces a flux change over
$t_{\mathrm{cross}}$, so it is $t_{\mathrm{cross}}$ --- set by the source
size --- that governs the observable variability. \checkmark

\textbf{(d)} At a quasar macro-image the stellar convergence is
$\kappab \sim 1$, meaning the projected separations of neighbouring stars are
comparable to their Einstein radii.  Their deflections then combine
non-linearly into an overlapping \emph{network} of caustics that covers the
source plane, rather than the isolated, well-separated point caustics of a
single Galactic lens.  The source is therefore microlensed continuously by the
ensemble, and its light curve reflects the caustic network it drifts across.
\checkmark
